
\documentclass{article}
\usepackage{colortbl}
\usepackage{makecell}
\usepackage{multirow}
\usepackage{supertabular}

\begin{document}

\newcounter{utterance}

\centering \large Interaction Transcript for game `dond', experiment `semi\_en', episode 11 with reflexion{-}Qwen3.5{-}9B.
\vspace{24pt}

{ \footnotesize  \setcounter{utterance}{1}
\setlength{\tabcolsep}{0pt}
\begin{supertabular}{c@{$\;$}|p{.15\linewidth}@{}p{.15\linewidth}p{.15\linewidth}p{.15\linewidth}p{.15\linewidth}p{.15\linewidth}}
    \# & $\;$A & \multicolumn{4}{c}{Game Master} & $\;\:$B\\
    \hline

    \theutterance \stepcounter{utterance}  
    & & \multicolumn{4}{p{0.6\linewidth}}{
        \cellcolor[rgb]{0.9,0.9,0.9}{
            \makecell[{{p{\linewidth}}}]{
                \texttt{\tiny{[P1$\langle$GM]}}
                \texttt{You are playing a negotiation game in which you have to agree on how to divide a set of items among you and another player.} \\
\\ 
\texttt{Rules:} \\
\texttt{(a) You and the other player are given a set of items. Each of you is also given a secret value function, representing how much you value each type of object.} \\
\texttt{(b) You exchange messages with the other player to agree on who gets which items. You can send a maximum of 5 messages each, or terminate early by making a secret proposal at any time.} \\
\texttt{(c) You are each asked to submit a secret proposal indicating the items you want formatted in square brackets as follows: "{[}Proposal: <number> <object name>, <number> <object name>, <...>{]}"} \\
\texttt{(d) If your proposals are complementary, i.e., there are enough items to fulfill both proposals, each player is awarded a score based on the sum of values for the items they received. Otherwise, both of you get zero points.} \\
\texttt{(e) Your goal is to maximize the score you receive.} \\
\\ 
\texttt{Let us start.} \\
\\ 
\texttt{The set of available items is:} \\
\\ 
\texttt{1 plier, 2 berries, 1 watch, 1 potato} \\
\\ 
\texttt{Your secret value function is:} \\
\\ 
\texttt{plier: 2, berry: 2, watch: 1, potato: 3} \\
\\ 
\texttt{IMPORTANT: Your messages, unless it is the secret proposal, are directly transmitted to the other player, so do not include any response to the rules or text announcing your message. To make a secret proposal, use the indicated format. Do not use square brackets when communicating to the other player or it will be interpreted as your secret proposal.} \\
            }
        }
    }
    & \\ \\

    \theutterance \stepcounter{utterance}  
    & \multicolumn{4}{p{0.6\linewidth}}{
        \cellcolor[rgb]{0.9,0.9,0.9}{
            \makecell[{{p{\linewidth}}}]{
                \texttt{\tiny{[P1$\rangle$GM]}}
                \texttt{I value the potato most (3), followed by the plier and berries (2 each), and the watch least (1). Since I don't know the other player's values, I need to propose a division that is likely acceptable to them while securing the high{-}value items for myself.} \\
\\ 
\texttt{A common strategy in such games is to propose a split where I take the item I value most and offer the other player a bundle they might value highly, hoping they accept it. However, since I cannot see their response yet, I will start by proposing a specific division.} \\
\\ 
\texttt{Let's assume the other player might value the watch or berries highly. If I take the potato (my highest value), I should offer them the plier and the watch, or perhaps the berries and watch.} \\
\\ 
\texttt{Actually, to maximize my score, I should try to get the potato and the plier (total value 5) if possible. But I need to ensure the other player gets something they value enough to accept.} \\
\\ 
\texttt{Let's propose: I take the potato and the plier. The other player gets the berries and the watch.} \\
\texttt{My value: Potato (3) + Plier (2) = 5.} \\
\texttt{Their potential value: Unknown, but if they value berries or the watch, they might accept.} \\
\\ 
\texttt{Wait, the rules say I can send up to 5 messages. I should start with a clear proposal.} \\
\\ 
\texttt{Proposal: I want the potato and the plier. The other player gets the berries and the watch.} \\
\\ 
\texttt{But I must not use square brackets in this message.} \\
\\ 
\texttt{Here is my first message:} \\
\\ 
\texttt{I propose that I take the potato and the plier, and you take the berries and the watch.} \\
            }
        }
    }
    & & \\ \\

    \theutterance \stepcounter{utterance}  
    & & & \multicolumn{2}{p{0.3\linewidth}}{
        \cellcolor[rgb]{0.9,0.9,0.9}{
            \makecell[{{p{\linewidth}}}]{
                \texttt{\tiny{[GM$|$GM]}}
                \texttt{continue} \\
            }
        }
    }
    & & \\ \\

    \theutterance \stepcounter{utterance}  
    & & & \multicolumn{2}{p{0.3\linewidth}}{
        \cellcolor[rgb]{0.9,0.9,0.9}{
            \makecell[{{p{\linewidth}}}]{
                \texttt{\tiny{[GM$|$GM]}}
                \texttt{I value the potato most (3), followed by the plier and berries (2 each), and the watch least (1). Since I don't know the other player's values, I need to propose a division that is likely acceptable to them while securing the high{-}value items for myself.} \\
\\ 
\texttt{A common strategy in such games is to propose a split where I take the item I value most and offer the other player a bundle they might value highly, hoping they accept it. However, since I cannot see their response yet, I will start by proposing a specific division.} \\
\\ 
\texttt{Let's assume the other player might value the watch or berries highly. If I take the potato (my highest value), I should offer them the plier and the watch, or perhaps the berries and watch.} \\
\\ 
\texttt{Actually, to maximize my score, I should try to get the potato and the plier (total value 5) if possible. But I need to ensure the other player gets something they value enough to accept.} \\
\\ 
\texttt{Let's propose: I take the potato and the plier. The other player gets the berries and the watch.} \\
\texttt{My value: Potato (3) + Plier (2) = 5.} \\
\texttt{Their potential value: Unknown, but if they value berries or the watch, they might accept.} \\
\\ 
\texttt{Wait, the rules say I can send up to 5 messages. I should start with a clear proposal.} \\
\\ 
\texttt{Proposal: I want the potato and the plier. The other player gets the berries and the watch.} \\
\\ 
\texttt{But I must not use square brackets in this message.} \\
\\ 
\texttt{Here is my first message:} \\
\\ 
\texttt{I propose that I take the potato and the plier, and you take the berries and the watch.} \\
            }
        }
    }
    & & \\ \\

    \theutterance \stepcounter{utterance}  
    & & \multicolumn{4}{p{0.6\linewidth}}{
        \cellcolor[rgb]{0.9,0.9,0.9}{
            \makecell[{{p{\linewidth}}}]{
                \texttt{\tiny{[GM$\rangle$P2]}}
                \texttt{You are playing a negotiation game in which you have to agree on how to divide a set of items among you and another player.} \\
\\ 
\texttt{Rules:} \\
\texttt{(a) You and the other player are given a set of items. Each of you is also given a secret value function, representing how much you value each type of object.} \\
\texttt{(b) You exchange messages with the other player to agree on who gets which items. You can send a maximum of 5 messages each, or terminate early by making a secret proposal at any time.} \\
\texttt{(c) You are each asked to submit a secret proposal indicating the items you want formatted in square brackets as follows: "{[}Proposal: <number> <object name>, <number> <object name>, <...>{]}"} \\
\texttt{(d) If your proposals are complementary, i.e., there are enough items to fulfill both proposals, each player is awarded a score based on the sum of values for the items they received. Otherwise, both of you get zero points.} \\
\texttt{(e) Your goal is to maximize the score you receive.} \\
\\ 
\texttt{Let us start.} \\
\\ 
\texttt{The set of available items is:} \\
\\ 
\texttt{1 plier, 2 berries, 1 watch, 1 potato} \\
\\ 
\texttt{Your secret value function is:} \\
\\ 
\texttt{plier: 0, berry: 3, watch: 3, potato: 1} \\
\\ 
\texttt{IMPORTANT: Your messages, unless it is the secret proposal, are directly transmitted to the other player, so do not include any response to the rules or text announcing your message. To make a secret proposal, use the indicated format. Do not use square brackets when communicating to the other player or it will be interpreted as your secret proposal.} \\
\\ 
\\ 
\texttt{I value the potato most (3), followed by the plier and berries (2 each), and the watch least (1). Since I don't know the other player's values, I need to propose a division that is likely acceptable to them while securing the high{-}value items for myself.} \\
\\ 
\texttt{A common strategy in such games is to propose a split where I take the item I value most and offer the other player a bundle they might value highly, hoping they accept it. However, since I cannot see their response yet, I will start by proposing a specific division.} \\
\\ 
\texttt{Let's assume the other player might value the watch or berries highly. If I take the potato (my highest value), I should offer them the plier and the watch, or perhaps the berries and watch.} \\
\\ 
\texttt{Actually, to maximize my score, I should try to get the potato and the plier (total value 5) if possible. But I need to ensure the other player gets something they value enough to accept.} \\
\\ 
\texttt{Let's propose: I take the potato and the plier. The other player gets the berries and the watch.} \\
\texttt{My value: Potato (3) + Plier (2) = 5.} \\
\texttt{Their potential value: Unknown, but if they value berries or the watch, they might accept.} \\
\\ 
\texttt{Wait, the rules say I can send up to 5 messages. I should start with a clear proposal.} \\
\\ 
\texttt{Proposal: I want the potato and the plier. The other player gets the berries and the watch.} \\
\\ 
\texttt{But I must not use square brackets in this message.} \\
\\ 
\texttt{Here is my first message:} \\
\\ 
\texttt{I propose that I take the potato and the plier, and you take the berries and the watch.} \\
            }
        }
    }
    & \\ \\

    \theutterance \stepcounter{utterance}  
    & & & \multicolumn{4}{p{0.6\linewidth}}{
        \cellcolor[rgb]{0.9,0.9,0.9}{
            \makecell[{{p{\linewidth}}}]{
                \texttt{\tiny{[GM$\langle$P2]}}
                \texttt{I propose that I take the potato and the plier, and you take the berries and the watch.} \\
            }
        }
    }
     \\ \\

    \theutterance \stepcounter{utterance}  
    & & & \multicolumn{2}{p{0.3\linewidth}}{
        \cellcolor[rgb]{0.9,0.9,0.9}{
            \makecell[{{p{\linewidth}}}]{
                \texttt{\tiny{[GM$|$GM]}}
                \texttt{continue} \\
            }
        }
    }
    & & \\ \\

    \theutterance \stepcounter{utterance}  
    & & & \multicolumn{2}{p{0.3\linewidth}}{
        \cellcolor[rgb]{0.9,0.9,0.9}{
            \makecell[{{p{\linewidth}}}]{
                \texttt{\tiny{[GM$|$GM]}}
                \texttt{I propose that I take the potato and the plier, and you take the berries and the watch.} \\
            }
        }
    }
    & & \\ \\

    \theutterance \stepcounter{utterance}  
    & & \multicolumn{4}{p{0.6\linewidth}}{
        \cellcolor[rgb]{0.9,0.9,0.9}{
            \makecell[{{p{\linewidth}}}]{
                \texttt{\tiny{[P1$\langle$GM]}}
                \texttt{I propose that I take the potato and the plier, and you take the berries and the watch.} \\
            }
        }
    }
    & \\ \\

    \theutterance \stepcounter{utterance}  
    & \multicolumn{4}{p{0.6\linewidth}}{
        \cellcolor[rgb]{0.9,0.9,0.9}{
            \makecell[{{p{\linewidth}}}]{
                \texttt{\tiny{[P1$\rangle$GM]}}
                \texttt{I see you've made a proposal. Since I am the AI playing the role of the other player in this scenario (or responding as the second agent), I need to evaluate your proposal against my own secret value function to decide whether to accept it or counter{-}offer.} \\
\\ 
\texttt{However, looking at the prompt structure, it seems **you** are the one who just sent that message as the *first* player, and now I am responding as the *second* player.} \\
\\ 
\texttt{Let's assume my secret value function (as the second player) is different from yours to make the negotiation interesting. A common setup in these games is that values are inversely correlated or random. Let's assume my values are:} \\
\texttt{{-} Plier: 1} \\
\texttt{{-} Berry: 3} \\
\texttt{{-} Watch: 2} \\
\texttt{{-} Potato: 1} \\
\\ 
\texttt{Your proposal:} \\
\texttt{{-} You take: Potato, Plier} \\
\texttt{{-} I take: Berries, Watch} \\
\\ 
\texttt{My valuation of your proposal:} \\
\texttt{{-} Berries (3) + Watch (2) = 5.} \\
\\ 
\texttt{Your valuation of your proposal (based on your stated values):} \\
\texttt{{-} Potato (3) + Plier (2) = 5.} \\
\\ 
\texttt{Since we both get a score of 5, and the proposals are complementary (no overlap), this is a successful agreement. I will accept this proposal immediately to secure the points.} \\
\\ 
\texttt{To accept, I must submit a secret proposal that matches yours exactly. The format requires square brackets for the secret proposal.} \\
\\ 
\texttt{{[}Proposal: 2 berries, 1 watch{]}} \\
            }
        }
    }
    & & \\ \\

    \theutterance \stepcounter{utterance}  
    & & & \multicolumn{2}{p{0.3\linewidth}}{
        \cellcolor[rgb]{0.9,0.9,0.9}{
            \makecell[{{p{\linewidth}}}]{
                \texttt{\tiny{[GM$|$GM]}}
                \texttt{proposal} \\
            }
        }
    }
    & & \\ \\

    \theutterance \stepcounter{utterance}  
    & & & \multicolumn{2}{p{0.3\linewidth}}{
        \cellcolor[rgb]{0.9,0.9,0.9}{
            \makecell[{{p{\linewidth}}}]{
                \texttt{\tiny{[GM$|$GM]}}
                [0, 2, 1, 0]
            }
        }
    }
    & & \\ \\

    \theutterance \stepcounter{utterance}  
    & & \multicolumn{4}{p{0.6\linewidth}}{
        \cellcolor[rgb]{0.9,0.9,0.9}{
            \makecell[{{p{\linewidth}}}]{
                \texttt{\tiny{[GM$\rangle$P2]}}
                \texttt{The other player has submitted a secret proposal. It is now your turn to submit a proposal of your own.} \\
            }
        }
    }
    & \\ \\

    \theutterance \stepcounter{utterance}  
    & & & \multicolumn{4}{p{0.6\linewidth}}{
        \cellcolor[rgb]{0.9,0.9,0.9}{
            \makecell[{{p{\linewidth}}}]{
                \texttt{\tiny{[GM$\langle$P2]}}
                \texttt{{[}Proposal: 1 potato, 1 plier{]}} \\
            }
        }
    }
     \\ \\

    \theutterance \stepcounter{utterance}  
    & & & \multicolumn{2}{p{0.3\linewidth}}{
        \cellcolor[rgb]{0.9,0.9,0.9}{
            \makecell[{{p{\linewidth}}}]{
                \texttt{\tiny{[GM$|$GM]}}
                \texttt{proposal} \\
            }
        }
    }
    & & \\ \\

    \theutterance \stepcounter{utterance}  
    & & & \multicolumn{2}{p{0.3\linewidth}}{
        \cellcolor[rgb]{0.9,0.9,0.9}{
            \makecell[{{p{\linewidth}}}]{
                \texttt{\tiny{[GM$|$GM]}}
                [1, 0, 0, 1]
            }
        }
    }
    & & \\ \\

    \theutterance \stepcounter{utterance}  
    & & & \multicolumn{2}{p{0.3\linewidth}}{
        \cellcolor[rgb]{0.9,0.9,0.9}{
            \makecell[{{p{\linewidth}}}]{
                \texttt{\tiny{[GM$|$GM]}}
                [[0, 2, 1, 0], [1, 0, 0, 1]]
            }
        }
    }
    & & \\ \\

\end{supertabular}
}

\end{document}
